\documentclass[11pt]{article}

% ---- Packages ----
\usepackage{amsmath,amssymb,amsthm}
\usepackage{mathtools}
\usepackage[hidelinks]{hyperref}
\usepackage{cleveref}
\usepackage{booktabs}
\usepackage{enumitem}
\usepackage[margin=1.15in]{geometry}
\usepackage{microtype}
\usepackage{xcolor}
\usepackage{tcolorbox}

% ---- Key-point boxes ----
\newtcolorbox{keybox}[1][]{
  colback=blue!4, colframe=blue!35,
  arc=3pt, boxrule=0.5pt,
  fonttitle=\bfseries, #1
}

% ---- Lean name formatting (line-breakable) ----
\ExplSyntaxOn
\tl_new:N \l__rfo_leanref_tl
\NewDocumentCommand \leanref { m }
  {
    \tl_set:Nn \l__rfo_leanref_tl { #1 }
    \regex_replace_all:nnN { \c{_} } { \c{_} \c{allowbreak} } \l__rfo_leanref_tl
    \regex_replace_all:nnN { \. } { . \c{allowbreak} } \l__rfo_leanref_tl
    \regex_replace_all:nnN { ([a-z])([A-Z]) } { \1 \c{allowbreak} \2 } \l__rfo_leanref_tl
    \regex_replace_all:nnN { ([0-9])([A-Z]) } { \1 \c{allowbreak} \2 } \l__rfo_leanref_tl
    \texttt{\small \tl_use:N \l__rfo_leanref_tl }
  }
\ExplSyntaxOff

% ---- Theorem environments ----
\theoremstyle{plain}
\newtheorem{theorem}{Theorem}[section]
\newtheorem{lemma}[theorem]{Lemma}
\newtheorem{corollary}[theorem]{Corollary}
\newtheorem{proposition}[theorem]{Proposition}

\theoremstyle{definition}
\newtheorem{definition}[theorem]{Definition}
\newtheorem{example}[theorem]{Example}

\theoremstyle{remark}
\newtheorem{remark}[theorem]{Remark}

% ---- Macros ----
\newcommand{\RFO}{\textsc{rfo}}
\newcommand{\RI}{\textsc{ri}}
\newcommand{\Lean}{\textsc{Lean\,4}}
\newcommand{\Mathlib}{\textsc{Mathlib}}
\newcommand{\ReflTransGen}{\operatorname{ReflTransGen}}
\newcommand{\PreservedBy}{\operatorname{PreservedBy}}
\newcommand{\ForwardClosed}{\operatorname{ForwardClosed}}
\newcommand{\reachableFrom}{\operatorname{reachableFrom}}
\newcommand{\unionGen}{\operatorname{unionGen}}
\newcommand{\obj}{\mathsf{obj}}
\newcommand{\mor}{\mathsf{mor}}
\newcommand{\represent}{\mathsf{represent}}

% ---- Title ----
\title{Reflective Fold Obstruction:\\
Internal Closure Does Not Exhaust Architectural Transition}
\author{Nova Spivack\\
Independent Researcher}
\date{2026}

\begin{document}
\maketitle

% ===========================================================================
% ABSTRACT
% ===========================================================================
\begin{abstract}
We prove a general theorem of \emph{fold obstruction}: if a predicate~$I$ is preserved
along the primitive steps of an internal relation~$r$, then $I$ is preserved along the
full reflexive--transitive closure~$r^{*}$, and any target failing~$I$ is not internally
reachable from a source satisfying~$I$.  In the hull form, when every seed in a
set~$S_0$ satisfies~$I$ and $I$ is forward-closed for~$r$, the entire closure hull
$\reachableFrom(r,S_0)$ satisfies~$I$; hence any $I$-violating target lies outside
that hull.

We instantiate this on the \emph{reflective slot calculus}, where sort structure (object
vs.\ morphism) is preserved by disciplined internal reflective steps, blocking any
internal route from the object slot to the morphism representer.  A bridge theorem
shows that every upstream \emph{representational system} inherits the reflective fold
obstruction.  Additional dynamic instances---including holonomy-phase and open--compact
witness obstructions with nontrivial step relations---are formalized in the artifact.

\emph{Semantic type obstruction} lifts the same hull mechanism to a preorder on
\emph{tags}: provably, \textbf{Turing-completeness does not imply semantic-type
completeness}.  A system may be maximally powerful in the standard computational sense and
yet be permanently bounded in which architectural modes (\emph{types of being}) it can
\emph{instantiate} along its \emph{own} primitive dynamics---distinct from what it can
\emph{simulate}.  We prove that \emph{typed simulations} pull specification-level type gaps
back to simulating systems, while a stricter \emph{realization} witness (a section with
backward step-lifting) is exactly the extra structure needed for a machine-checked
\emph{bidirectional} agreement of type reachability between a model and a pulled-back
typing.

The entire development is machine-checked in \Lean{}/\Mathlib{} with zero \texttt{sorry}
and zero \texttt{admit} across all shipped modules.
\end{abstract}

% ===========================================================================
% 1. THE PROBLEM
% ===========================================================================
\section{The problem: internal closure versus architectural transition}\label{sec:problem}

Consider a system whose states evolve along an \emph{internal} step relation
$r : \alpha \to \alpha \to \mathrm{Prop}$.  Reflexive--transitive closure~$r^{*}$
captures everything the system can reach by iterating, composing, and idling along~$r$.
A natural question is:

\begin{quote}
\emph{Does internal closure exhaust architectural transition?  That is, can every state
reachable ``in principle'' be obtained by iterating~$r$?}
\end{quote}

This paper answers: \textbf{no}, in a precise and general sense.

We prove that whenever a predicate~$I$ is preserved by each primitive~$r$-step, the
predicate lifts to the entire closure~$r^{*}$.  Any target violating~$I$ is therefore
\emph{not} internally reachable from a source satisfying~$I$.  We call such a blocked
transition a \emph{fold}: a structural passage that the internal calculus cannot
generate.

\begin{keybox}[title={A fold is not an iterate}]
The central claim of this paper: any architecture generated solely by internal
reachability remains confined to its invariant class.  Mismatch across that boundary
certifies that the target requires a \emph{fold}---a passage not generated by~$r^{*}$.
Internal closure and architecture-changing transition are structurally different
operations.
\end{keybox}

\noindent The claim is not that the unreachable target is impossible in all senses, but
that any route generated only by~$r$ and preserving~$I$ cannot reach it.  Internal
closure and architecture-changing transition must therefore be distinguished.

\medskip
\noindent
\textbf{A second orthogonal ceiling (preview).}
Classical computability answers how much a Turing-complete model can \emph{compute};
this paper's fold machinery answers how far a fixed architecture can move \emph{within
itself} under a chosen primitive dynamics.  These are independent: internal iteration may
preserve invariants that cap \emph{semantic types}---modes of self-relation or
adjudication the system \emph{instantiates}---even when no computability-theoretic barrier
blocks \emph{simulating} deeper behaviour externally.  \Cref{sec:semantic-type} makes this
split precise at the level of machine-checked \emph{type reachability} and simulation,
and \Cref{sec:conclusion} summarizes the stakes.

% ===========================================================================
% 2. DEFINITIONS
% ===========================================================================
\section{Definitions}\label{sec:defs}

We work over a type~$\alpha$ in a universe~$u$.  All definitions and theorems in this
section have exact counterparts in the \Lean{} artifact.

\begin{definition}[Primitive internal relation]\label{def:relation}
A \emph{primitive internal relation} on~$\alpha$ is a binary relation
$r : \alpha \to \alpha \to \mathrm{Prop}$.  We make no assumptions on~$r$ (it need
not be reflexive, transitive, or decidable).
\end{definition}

\begin{definition}[Forward closure / preserved invariant]\label{def:forwardclosed}
A predicate $I : \alpha \to \mathrm{Prop}$ is \emph{forward-closed} for~$r$
(equivalently, \emph{preserved by}~$r$) when
\[
  \ForwardClosed(r,I) \;\coloneqq\;
  \forall\, a\, b,\;\; r(a,b) \;\Longrightarrow\; I(a) \;\Longrightarrow\; I(b).
\]
We write $\PreservedBy(r,I)$ as a synonym.
\end{definition}

\begin{definition}[Reflexive--transitive closure]\label{def:refltransgen}
$\ReflTransGen(r)$ is the reflexive--transitive closure of~$r$, the standard
Mathlib~\cite{mathlib} \texttt{Relation.ReflTransGen}:
\begin{itemize}[nosep]
  \item $\ReflTransGen(r,a,a)$ for every~$a$ \,(\emph{refl});
  \item if $\ReflTransGen(r,a,b)$ and $r(b,c)$, then $\ReflTransGen(r,a,c)$
    \,(\emph{tail}).
\end{itemize}
\end{definition}

\begin{definition}[Closure hull]\label{def:hull}
Given a seed set $S_0 \subseteq \alpha$, the \emph{closure hull} is
\[
  \reachableFrom(r, S_0) \;=\;
  \bigl\{\, y \in \alpha \;\bigm|\;
  \exists\, x \in S_0,\; \ReflTransGen(r,x,y) \,\bigr\}.
\]
\end{definition}

% ===========================================================================
% 3. MAIN THEOREM
% ===========================================================================
\section{The main theorem}\label{sec:main}

The theorem is completely general in the carrier type~$\alpha$ and primitive
relation~$r$.  No categorical, topological, or computational assumptions are required
beyond preservedness of the chosen invariant~$I$.

\subsection*{Theorem A: primitive-to-closure preservation}

\begin{theorem}[Closure preservation]\label{thm:closure-preservation}
Let $r : \alpha \to \alpha \to \mathrm{Prop}$ and $I : \alpha \to \mathrm{Prop}$ with
$\ForwardClosed(r,I)$.  Then for all $a,b : \alpha$,
\[
  \ReflTransGen(r,a,b) \;\Longrightarrow\; I(a) \;\Longrightarrow\; I(b).
\]
\end{theorem}

\begin{proof}
By induction on the derivation of $\ReflTransGen(r,a,b)$.  The \emph{refl} case is
immediate.  For the \emph{tail} case, given $\ReflTransGen(r,a,b')$ and $r(b',b)$, the
inductive hypothesis yields $I(b')$, and forward closure gives $I(b)$.
\end{proof}

\smallskip
\noindent\textbf{Lean:}
\leanref{internal\_reachability\_preserves\_invariant}.

\subsection*{Theorem B: pairwise fold obstruction}

\begin{theorem}[Fold obstruction]\label{thm:fold-obstruction}
If $\PreservedBy(r,I)$, $I(S)$, and $\lnot I(T)$, then
\[
  \lnot\, \ReflTransGen(r, S, T).
\]
\end{theorem}

\begin{proof}
Suppose for contradiction that $\ReflTransGen(r,S,T)$.
By \Cref{thm:closure-preservation}, $I(S)$ implies $I(T)$, contradicting $\lnot I(T)$.
\end{proof}

\smallskip
\noindent\textbf{Lean:}
\leanref{fold\_obstruction\_of\_invariant\_mismatch}.

\begin{keybox}[title={The master obstruction}]
Any architecture generated solely by internal reachability remains confined to its
invariant class; mismatch across that boundary certifies that the target requires a
fold.  The obstacle is structural: no amount of iteration, composition, or closure
under~$r$ can cross the boundary drawn by a preserved~$I$.
\end{keybox}

\begin{remark}\label{rem:interpretation}
\Cref{thm:fold-obstruction} does not say that $T$ is ``impossible.''  It says that any
path generated exclusively by~$r$ from a source satisfying~$I$ cannot reach a target
violating~$I$.  Hence internal closure and invariant-breaking transition must be
distinguished: the latter requires a \emph{fold}---a structural passage not in~$r^{*}$.
\end{remark}

\subsection*{Theorem C: hull exclusion}

\begin{theorem}[Hull preservation]\label{thm:hull-preservation}
If $\PreservedBy(r,I)$ and $\forall s \in S_0,\; I(s)$, then
\[
  \forall\, x \in \reachableFrom(r, S_0),\; I(x).
\]
\end{theorem}

\begin{proof}
Let $x \in \reachableFrom(r,S_0)$, so there exists $s \in S_0$ with
$\ReflTransGen(r,s,x)$.  Since $I(s)$ holds and $I$ is forward-closed,
\Cref{thm:closure-preservation} gives $I(x)$.
\end{proof}

\begin{corollary}[Hull exclusion]\label{cor:hull-exclusion}
Under the hypotheses of \Cref{thm:hull-preservation}, if $\lnot I(T)$, then
$T \notin \reachableFrom(r,S_0)$.
\end{corollary}

\smallskip
\noindent\textbf{Lean:}
\leanref{reachableFrom\_hull\_preserves\_invariant},\\
\leanref{not\_mem\_reachableFrom\_of\_preserved\_invariant\_mismatch}.

% ===========================================================================
% 4. REFLECTIVE INSTANTIATION
% ===========================================================================
\section{Reflective instantiation: sort obstruction}\label{sec:reflective}

The abstract theorem is not vacuous.  We instantiate it on a canonical reflective
architecture, demonstrating a genuine dynamical obstruction at the \emph{sort} level.

\begin{definition}[Reflective system]\label{def:reflective-system}
A \emph{reflective system} $R$ consists of a category~$\mathcal{C}$, a distinguished
object $A \in \mathcal{C}$, and a self-representation endomorphism
$\represent : A \to A$.  This is the minimal categorical setting in which ``iterate the
self-representation internally'' becomes mathematically meaningful: the endomorphism
monoid $\mathrm{End}(A)$ supplies the iteration, and the category supplies the ambient
structure~\cite{lawvere1969,yanofsky2003}.
\end{definition}

\begin{definition}[Ontological slot]\label{def:slot}
The \emph{ontological slot} type on~$R$ is the disjoint union
\[
  \mathrm{Slot}(R) \;=\; \obj(\mathcal{C}) \;+\; \mor(\mathrm{End}(A)),
\]
tagging each element as either an \emph{object tag} $\obj(c)$ for $c \in \mathcal{C}$,
or a \emph{morphism tag} $\mor(f)$ for $f : A \to A$.  These constructors are injective
and disjoint: $\obj(c) \neq \mor(f)$ for all $c, f$.
\end{definition}

\begin{definition}[Reflective slot step]\label{def:slotstep}
The primitive internal relation $\mathsf{reflectiveSlotStep}(R)$ on $\mathrm{Slot}(R)$
is defined by:
\begin{itemize}[nosep]
  \item $s = t$ \,(\emph{idle}); or
  \item $s = \mor(f)$, $t = \mor(g)$, and $g = f \circ \represent$ \,(\emph{mor advance}).
\end{itemize}
Object slots cannot be stepped into morphism slots: the only moves from $\obj(c)$ are
reflexive.
\end{definition}

\begin{definition}[Object-branch predicate]\label{def:isobjslot}
$\mathrm{IsObj}(s) \coloneqq \exists\, c,\; s = \obj(c)$.
\end{definition}

\begin{lemma}[Sort preservation]\label{lem:sort-preservation}
$\ForwardClosed(\mathsf{reflectiveSlotStep}(R),\, \mathrm{IsObj})$.
\end{lemma}

\begin{proof}
If $s$ is an object tag $\obj(c)$ and $\mathsf{reflectiveSlotStep}(R,s,t)$, then either
$t = s$ (idle), or $s$ would need to be a morphism tag (impossible for $\obj(c)$).
Hence $t$ is still an object tag.
\end{proof}

\begin{lemma}\label{lem:represent-not-obj}
$\lnot\, \mathrm{IsObj}\bigl(\mor(\represent)\bigr)$.
\end{lemma}

\begin{proof}
If $\mor(\represent) = \obj(c)$ for some~$c$, this contradicts the constructor
disjointness of the slot type.
\end{proof}

\begin{theorem}[Reflective fold obstruction --- sort mismatch]\label{thm:reflective-fold}
For any reflective system~$R$,
\[
  \lnot\, \ReflTransGen\bigl(\mathsf{reflectiveSlotStep}(R),\;
  \obj(A),\; \mor(\represent)\bigr).
\]
\end{theorem}

\begin{proof}
By \Cref{thm:fold-obstruction} with $I \coloneqq \mathrm{IsObj}$, using
\Cref{lem:sort-preservation} (preservation), $\langle A, \mathrm{rfl}\rangle$
(source satisfies $I$), and \Cref{lem:represent-not-obj} (target violates $I$).
\end{proof}

\smallskip
\noindent\textbf{Lean:}
\leanref{ReflectiveFold.reflective\_fold\_obstruction\_slot\_mismatch}.

\begin{keybox}[title={Sort structure is a fold boundary}]
In a reflective system, the passage from the object slot $\obj(A)$ to the morphism
representer $\mor(\represent)$ is \textbf{not an iterate}: no chain of internal
reflective steps can cross the sort boundary.  The self-representation endomorphism
generates unbounded depth within the morphism branch, but it cannot bridge the
structural gap between object and morphism.  This is the precise sense in which
\textbf{a fold is not an iterate}.
\end{keybox}

\begin{remark}[Rich reflective calculus]\label{rem:rich-calculus}
The obstruction survives under the \emph{rich reflective calculus}
$\mathsf{reflectiveCalcStep}(R)$, which extends the primitive step by allowing
mor-branch jumps $f \mapsto f \circ \represent^n$ for \emph{any} positive iterate
$n \geq 1$.  Sort preservation holds by the same argument: object tags admit no
morphism-branch transitions.  Under the hypothesis $\mathrm{IterInjective}(R)$
(distinct natural-number exponents give distinct endomorphism powers), the rich
calculus is \emph{strictly} stronger than the primitive slot step.
\end{remark}

% ===========================================================================
% 5. DYNAMIC INSTANCES
% ===========================================================================
\section{A dynamic instance: holonomy-phase obstruction}\label{sec:instances}

The reflective slot obstruction of \Cref{sec:reflective} is the primary instance.  To
show the abstract theorem is not limited to sort-level reasoning, we highlight one
further instance with genuinely nontrivial dynamics.

\begin{example}[Holonomy-phase obstruction]\label{ex:holonomy}
The carrier is $\mathrm{HolonomyState}$, a gated phase counter.  The internal relation
$\mathsf{holonomyPhaseStep}$ advances the phase under a nontrivial gating condition---it
is \emph{not} mere equality.  The invariant $\mathrm{HolonomyTrivialBounded}$
(trivial holonomy tag with bounded phase) is forward-closed for this step.  A target
state with twist holonomy lies outside the closure hull of trivially-tagged seeds.
Structurally, this means that a system whose holonomy phase evolves under the gated
dynamics \emph{cannot} reach a twisted configuration from a trivial one by internal
iteration alone: the twist requires a fold.
\end{example}

\noindent The machine-checked artifact formalizes additional instances: boundary-type,
connected-boundary, orientability-parity, holonomy-tag, and open--compact witness
obstructions.  Each is a one-line corollary of \Cref{thm:fold-obstruction} with its own
carrier, invariant, and step relation.  These are catalogued in the artifact's theorem
inventory.

% ===========================================================================
% 6. RELATION TO REPRESENTATIONAL INCOMPLETENESS
% ===========================================================================
\section{Relation to representational incompleteness}\label{sec:ri}

Reflective fold obstruction (\RFO) and representational incompleteness
(\RI)~\cite{spivack2026ri} are complementary programs.

\RI{} proves that no finitely-typed system can universally parametrize its own function
space: the diagonal / Lawvere barrier~\cite{lawvere1969} blocks total self-capture.  The
obstacle is a \emph{cardinality}-flavored impossibility (surjectivity failure of curry).

\RFO{} proves that certain architectural transitions---those that would violate a
preserved invariant---cannot be realized by internal closure.  The obstacle is a
\emph{transport}-flavored impossibility (invariant mismatch along reachability).

Atop that dynamical core, \Cref{sec:semantic-type} packages a \emph{semantic typing}
layer: tags, instantiation, and type reachability.  It answers a question orthogonal to
\RI's diagonal barrier---how internal dynamics bounds \emph{architectural kinds}---while
still machine-checkable without importing a full Rice-style index-set formalization.

The two programs share reflective-system infrastructure but own different theorems.
Every upstream \RI{} host inherits the \RFO{} sort obstruction through a proved
bridge:

\begin{theorem}[RI bridge]\label{thm:ri-bridge}
For any $\mathrm{RepresentationalSystem}$~$R$ (in the sense of~\RI{}), the induced
$\mathrm{ReflectiveSystem}$ satisfies
\[
  \lnot\, \ReflTransGen\bigl(\mathsf{reflectiveSlotStep},\;
  \obj(A),\; \mor(\represent)\bigr).
\]
\end{theorem}

\noindent\textbf{Lean:}
\leanref{representational\_incompleteness\_implies\_reflective\_fold\_obstruction}.

\begin{remark}
The bridge adds \emph{no new axioms}: it is definitional alignment plus the generic
reflective slot theorem.  The hypothesis $\mathrm{IterInjective}$ is required only for
the tower-unboundedness channel (a separate, complementary obstruction), not for the
fold statement.
\end{remark}

% ===========================================================================
% 7. ROBUSTNESS
% ===========================================================================
\section{Robustness of the obstruction}\label{sec:robustness}

The fold pattern is not brittle.  Three meta-stability results ensure that the
obstruction survives routine extensions of the primitive relation.

\begin{proposition}[Subrelation persistence]\label{prop:subrelation}
If $r \subseteq r'$ as graphs and $\ForwardClosed(r', I)$, then $\ForwardClosed(r, I)$.
The fold obstruction for~$r$ follows from preservation on a single larger
relation~$r'$, without a second induction.
\end{proposition}

\begin{proposition}[Generator-wise preservation]\label{prop:generator}
Given an indexed family $\{r_i\}_{i \in \iota}$, define the generated one-step relation
$\unionGen(\{r_i\})(x,y) \coloneqq \exists\, i,\; r_i(x,y)$.  If
$\ForwardClosed(r_i, I)$ for every~$i$, then $\ForwardClosed(\unionGen(\{r_i\}), I)$,
and the fold obstruction holds for the generated calculus.
\end{proposition}

\noindent Additionally, for any indexed family of admissible reflective generators,
each preserving sort separation, the union relation inherits the fold obstruction.  The
single-step rich calculus and its $\ReflTransGen$-hull are shown to lie inside the
canonical generator-union hull, so the generated obstruction subsumes the hand-written
one.

\begin{keybox}[title={The obstruction is stable}]
Strengthening the internal calculus---adding new primitive steps, taking unions of
generator families, or passing to richer reflective dynamics---does not escape the
obstruction.  As long as the enlarged relation still preserves the invariant, the fold
boundary holds.  This makes the theorem a genuine \emph{obstruction calculus}, not a
fragile artifact of a specific step system.
\end{keybox}

% ===========================================================================
% 8. SIGNIFICANCE
% ===========================================================================
\section{Significance}\label{sec:significance}

\Cref{thm:fold-obstruction,thm:reflective-fold} together establish a clean boundary
between what internal closure can generate and what it cannot.

\begin{enumerate}[label=(\arabic*),nosep,leftmargin=2em]
  \item \textbf{The theorem separates internal closure from architecture-changing
    transition.}  Any route generated only by the chosen internal relation~$r$, from a
    source satisfying~$I$, cannot reach a target violating~$I$.  Passages that cross an
    invariant boundary require a \emph{fold}---a transition not generated by~$r$.

  \item \textbf{The reflective instance shows the theorem is nonvacuous in a canonical
    self-referential setting.}  Sort structure is preserved under disciplined internal
    reflective dynamics, so the object-to-representer passage is genuinely blocked.

  \item \textbf{The artifact proves this is a machine-checked obstruction calculus, not
    an informal schema.}  The entire development, from abstract definitions through
    concrete instances, compiles with zero proof holes.

  \item \textbf{Semantic types vs.\ Turing completeness (\Cref{sec:semantic-type}).}
    The library proves that the preorder of \emph{semantic-type reachability} is
    genuinely nontrivial and that \emph{internal} type gaps persist regardless of how
    much ``raw computation'' the host could in principle support: simulation of a deeper
    configuration is not instantiation of its type along the host's primitive step.
    Typed simulations preserve forward reachability and pull \emph{obstructions} back;
    a proved \emph{section} criterion identifies when this becomes an \emph{iff}---the
    formal edge between mere simulation and \emph{realization} at the typing layer.
\end{enumerate}

\subsection*{Novelty}

We do not claim that preserved-invariant reasoning or reflexive--transitive
closure~\cite{barendregt1984,mathlib} are individually new.  A central novelty of this
paper is the formal realization of the principle that \emph{a fold is not an iterate}:
internally generated reachability, even under closure and enriched calculi, remains
confined to invariant-preserving classes, so invariant-breaking architectural transitions
cannot be obtained by iteration alone.  The contribution is the
\textbf{machine-checked integration} in which this principle is proved, instantiated on
reflective architectures, extended to generated calculi and admissible families, and
bridged to upstream representational incompleteness---all with zero proof holes.

\begin{keybox}[title={The novelty is the formal realization}]
Individual ingredients are standard.  What is new is a machine-checked integration in
which preserved invariants lift from primitive steps to reflexive--transitive closure,
yielding pairwise obstruction, hull exclusion, reflective sort obstruction,
generated-calculus stability, and bridge compatibility in one theorem stack---so that
internal iteration alone is provably insufficient for invariant-breaking architectural
transition.  The same stack now includes a \emph{semantic-type} layer with explicit
theorems on pullbacks, simulation, complementary-tag gaps under one-sided preservation,
and seed-local hull exclusion---articulating limits on Turing-complete hosts that are not
the classical halting/Rice storyline, but dynamical obstruction at the level of what a
system \emph{is} under its own transition relation.
\end{keybox}

% ===========================================================================
% 9. SEMANTIC TYPE OBSTRUCTION (SPEC\_020)
% ===========================================================================
\section{Semantic type obstruction}\label{sec:semantic-type}

\subsection*{Turing completeness and an orthogonal semantic ceiling}

Classical computability fixes a system's \emph{computational} ceiling: under standard
notions of algorithm, a Turing-complete model can compute every partial recursive function
and no others.  That theorem classifies \emph{external} functional behaviour over standard
presentations of algorithms.  It does \emph{not} classify what \emph{semantic types} a
system can \emph{instantiate}---modes of self-relation, self-modeling, or adjudication---as
matters of architectural \emph{being} along its \emph{own} internal transition relation.

\medskip
The machine-checked layer below shows these dimensions are \textbf{orthogonal}.  A host may
be Turing-complete (hence able to simulate rich counterfactual dynamics) while its
primitive calculus preserves predicates that make certain tag pairs \emph{internally}
unreachable forever: no sequence of admissible internal steps joins a source tag to a
target tag, even though the host could \emph{simulate} a machine that appears to realize
the target tag externally.  What is proved is not a replacement for Rice or the halting
problem, but a \emph{different} impossibility template: \textbf{invariant preservation
along $r$ closes the hull}, at the state level and then at the tag level.

\subsection*{Semantic typing}

Fix a carrier of \emph{systems}, a discrete set of \emph{tags} (semantic types), a binary
\emph{primitive step} relation, and a predicate stating when a system \emph{instantiates}
a tag.  If every primitive step preserves each tag's extension, then any reflexive--transitive
internal route from a system instantiating tag~$T$ stays inside that tag's invariant class.
Thus two tags separated by an invariant mismatch are not connected by internal iteration;
this is the same hull mechanism as \Cref{thm:fold-obstruction}, lifted to a preorder on tags.

\subsection*{Non-triviality}

The \emph{type reachability} preorder on tags is not total: concrete witnesses exhibit
pairs $(T,T')$ with no connecting internal chain.  One canonical family uses
\emph{self-modeling depth} labels on~$\mathbb{N}$ with identity dynamics: depth is
preserved, so tag~\texttt{depth(0)} cannot reach \texttt{depth(1)}---a crisp toy model of
``type-bounded'' introspection.  A second family abstracts \emph{adjudication class}
(total-effective vs.\ weaker behaviour): kind is preserved along steps, so the total tag
cannot reach the non-total tag without a fold.  A third instance reframes the reflective
\emph{sort} gap: the object-branch tag cannot reach the morphism-carrier tag under
\texttt{reflectiveSlotStep}, so the diagonal witness \texttt{mor(represent)} is unreachable
from the parametric object branch---in parallel with the RI bridge (\Cref{sec:ri}).

\subsection*{Simulation, pullback, and realization}

To compare architectures, the library formalizes \emph{typed primitive simulations}: a map
$\pi$ from a domain state space to states of a fixed typing, carrying the domain's primitive
relation into the typing's primitive steps.  Pulling tags back along $\pi$ yields a typing
on the domain (\texttt{pullbackAlongSimulation}).  \textbf{Forward simulation} is always
available: type reachability in the pullback implies type reachability in the codomain
(via a homomorphism lift of \texttt{ReflTransGen}).

\medskip
\noindent
\textbf{The converse is conditional---and this is the formal simulation/realization split.}
In general one cannot expect backward reachability in the pullback from the codomain alone:
simulation need not exhibit all trajectories of the simulated system in the host's native
dynamics.  The artifact isolates a minimal witness that repairs the gap: a \emph{section}
$\sigma$ with $\pi \circ \sigma = \mathrm{id}$ that \emph{lifts every primitive step backward}
(\texttt{TypedPrimitiveSimulationSection}).  Under exactly this data, \textbf{type
reachability is equivalent} between the codomain typing and the pullback, and the same
holds for packaged \texttt{SemanticTypeObstruction} (\texttt{typeReachable\_pullback\_iff\_of\_section}).
Thus ``realizing'' the foreign dynamics inside the host's step relation is a strictly
stronger claim than forward simulation, and the library states the strengthening in
provable \texttt{iff} form rather than by axiom.

\subsection*{Complementary tags and seed-local hulls (weaker hypotheses)}

Many applications single out one invariant $P$ whose forward-closure is cheap, without
assuming a fully bipartite typing where $\neg P$ is also step-closed.  Two results capture
this cleanly: \textbf{complementary-tag gaps} (\texttt{typeGap\_of\_complementary\_tags})
propagate from two tags that classify $P$ and $\neg P$ respectively, requiring only that
$P$ itself be forward-closed along primitive steps; and \textbf{seed-local hull exclusion}
(\texttt{false\_tag\_not\_mem\_tagsOnHull\_of\_preserved\_seed}) shows that a tag for
$\neg P$ never appears on the primitive hull from $P$-seeds under the same one-sided
hypothesis.  These lemmas are obstruction-theoretic footnotes to the master hull theorem,
but they match how practitioners reason about reachable components from real initial
conditions.

\subsection*{Further corollaries and scope}

The library also includes: \emph{primitive extension} monotonicity; \emph{tags-on-hull}
from seeds; \emph{pair-tag} conjunction obstruction; a \emph{bipartition} bridge to
\texttt{FoldObstructionBundle} when both $P$ and $\neg P$ are forward-closed; an infinite
\emph{pairwise antichain} of depth labels; and an explicit computability \emph{scope note}
contrasting Rice-style index phenomena with this \emph{dynamical} formulation (no
smuggled formalization of r.e.\ degrees is required to use the typing theorems).

\subsection*{Lean anchors}

The development lives under \texttt{ReflectiveFoldObstruction/SemanticType/} with stable
import surface \texttt{Core/SemanticTypeObstruction}.  Core names include
\texttt{semanticType\_preorder\_nontrivial},
\texttt{selfModelDepth\_obstruction},
\texttt{adjudication\_semantic\_obstruction},
\texttt{RI\_semantic\_type\_mismatch}; simulation/realization:
\texttt{typeGap\_simulation\_pullback},
\texttt{typeReachable\_pullback\_iff\_of\_section},
\texttt{SemanticTypeObstruction\_pullback\_iff\_of\_section}; complementary/hull:
\texttt{typeGap\_of\_complementary\_tags},
\texttt{false\_tag\_not\_mem\_tagsOnHull\_of\_preserved\_seed}.
See \Cref{tab:anchors}.

% ===========================================================================
% 10. MACHINE-CHECKED ARTIFACT
% ===========================================================================
\section{Machine-checked artifact}\label{sec:artifact}

\begin{center}
\begin{tabular}{@{}ll@{}}
\toprule
\textbf{Component} & \textbf{Value} \\
\midrule
Proof assistant & \Lean{} \\
Toolchain & \texttt{leanprover/lean4:v4.29.0-rc6} \\
Library & \Mathlib{} v4.29.0-rc6 \\
Build command & \texttt{lake build ReflectiveFoldObstruction} \\
\texttt{sorry} count & 0 \\
\texttt{admit} count & 0 \\
\bottomrule
\end{tabular}
\end{center}

\subsection*{Theorem anchors}

\Cref{tab:anchors} lists the principal Lean names.

\begin{table}[h]
\centering
\small
\begin{tabular}{@{}p{0.33\textwidth}p{0.60\textwidth}@{}}
\toprule
\textbf{Role} & \textbf{Lean name} \\
\midrule
Closure preservation (\Cref{thm:closure-preservation})
  & \leanref{internal\_reachability\_preserves\_invariant} \\[3pt]
Fold obstruction (\Cref{thm:fold-obstruction})
  & \leanref{fold\_obstruction\_of\_invariant\_mismatch} \\[3pt]
Hull preservation (\Cref{thm:hull-preservation})
  & \leanref{reachableFrom\_hull\_preserves\_invariant} \\[3pt]
Hull exclusion (\Cref{cor:hull-exclusion})
  & \leanref{not\_mem\_reachableFrom\_of\_preserved\_invariant\_mismatch} \\[3pt]
Sort preservation (\Cref{lem:sort-preservation})
  & \leanref{reflective\_step\_preserves\_sort\_separation} \\[3pt]
Reflective fold (\Cref{thm:reflective-fold})
  & \leanref{ReflectiveFold.reflective\_fold\_obstruction\_slot\_mismatch} \\[3pt]
Rich calculus fold (\Cref{rem:rich-calculus})
  & \leanref{reflectiveCalc\_fold\_obstruction\_slot\_mismatch} \\[3pt]
\RI{} bridge (\Cref{thm:ri-bridge})
  & \leanref{representational\_incompleteness\_implies\_reflective\_fold\_obstruction} \\[3pt]
Subrelation persistence (\Cref{prop:subrelation})
  & \leanref{preserved\_under\_relation\_extension} \\[3pt]
Generated-calculus fold (\Cref{prop:generator})
  & \leanref{generated\_calculus\_fold\_obstruction} \\[3pt]
Semantic-type non-totality (\Cref{sec:semantic-type})
  & \leanref{semanticType\_preorder\_nontrivial} \\[3pt]
RI semantic mismatch (\Cref{sec:semantic-type})
  & \leanref{RI\_semantic\_type\_mismatch} \\[3pt]
Relation map / simulation (\Cref{sec:semantic-type})
  & \leanref{reflTransGen\_map}, \leanref{typeGap\_simulation\_pullback},
    \leanref{typeReachable\_pullback\_iff\_of\_section},
    \leanref{typeGap\_of\_complementary\_tags},
    \leanref{false\_tag\_not\_mem\_tagsOnHull\_of\_preserved\_seed} \\
\bottomrule
\end{tabular}
\caption{Principal machine-checked theorem anchors.}\label{tab:anchors}
\end{table}

\subsection*{Reproduction}

\begin{verbatim}
  cd reflective-fold-obstruction-lean
  lake update && lake exe cache get
  lake build ReflectiveFoldObstruction
\end{verbatim}

\noindent
The full module tree and formalization map are given in \Cref{app:modules}.

% ===========================================================================
% 11. CONCLUSION
% ===========================================================================
\section{Conclusion}\label{sec:conclusion}

This paper's core contribution is a general, machine-checked \emph{fold obstruction}:
preserved invariants propagate along reflexive--transitive closure, so any target that
violates the invariant lies outside internal reachability from compliant sources.  The
principle is instantiated on reflective sort dynamics, extended to rich calculi and
generator unions, and bridged to representational incompleteness without new axioms.

\medskip
The \emph{semantic type} layer (\Cref{sec:semantic-type}) pushes the same idea to tags
and makes vivid a limit that is easy to conflate with classical computability but is
logically distinct: \textbf{Turing-completeness does not entail semantic-type
completeness}.  A system may compute everything a textbook universal machine computes,
while its native step relation preserves architectural predicates that \emph{permanently}
exclude certain semantic kinds from the internally reachable class.  Simulation can mask
this from outside observers; what the formalization rewards is careful separation between
\emph{what can be forward-simulated} and \emph{what is realized along the host's own
primitive transitions}---with a proved \texttt{iff} characterizing the latter in terms of
backward-lifting sections.

\medskip
Together, the development supports a unified moral for reflective architectures and their
machine-checked obstructions: \textbf{internal closure is not architectural
universality}.  Iteration along a disciplined dynamics respects invariants; crossing an
invariant boundary is a fold, not an iterate---whether the invariant is sort structure,
holonomy data, boundary connectivity, semantic adjudication class, or the instantiation
of a finer semantic tag.

% ===========================================================================
% REFERENCES
% ===========================================================================
\begin{thebibliography}{99}

\bibitem{lawvere1969}
F.~W.~Lawvere.
\newblock Diagonal arguments and cartesian closed categories.
\newblock In \emph{Category Theory, Homology Theory and Their Applications~II},
  Lecture Notes in Mathematics, vol.~92, pp.~134--145. Springer, 1969.

\bibitem{yanofsky2003}
N.~S.~Yanofsky.
\newblock A universal approach to self-referential paradoxes, incompleteness and
  fixed points.
\newblock \emph{Bulletin of Symbolic Logic}, 9(3):362--386, 2003.

\bibitem{barendregt1984}
H.~P.~Barendregt.
\newblock \emph{The Lambda Calculus: Its Syntax and Semantics}.
\newblock Studies in Logic and the Foundations of Mathematics, vol.~103.
  North-Holland, revised edition, 1984.

\bibitem{mathlib}
The Mathlib Community.
\newblock \emph{Mathlib}: the Lean mathematical library.
\newblock \url{https://github.com/leanprover-community/mathlib4}, 2020--2026.

\bibitem{lean4}
L.~de~Moura and S.~Ullrich.
\newblock The \emph{Lean~4} theorem prover and programming language.
\newblock In \emph{CADE-28}, Lecture Notes in Computer Science, vol.~12699,
  pp.~625--635. Springer, 2021.

\bibitem{spivack2026ri}
N.~Spivack.
\newblock Representational incompleteness: no self-modeling system can capture its own
  diagonal.
\newblock Preprint, 2026.
\newblock Machine-checked in Lean~4/Mathlib.

\end{thebibliography}

% ===========================================================================
% APPENDIX
% ===========================================================================
\appendix

\section{Module map}\label{app:modules}

\small
\begin{center}
\begin{tabular}{@{}p{0.45\textwidth}p{0.49\textwidth}@{}}
\toprule
\textbf{Module path} & \textbf{Role} \\
\midrule
\texttt{Core/Basic} & \texttt{ReflectiveSystem}, \texttt{IterInjective}, iterates, monoid laws \\
\texttt{Core/Slots} & \texttt{OntologicalSlot}, sort disjointness, \texttt{ReflectiveSlot} \\
\texttt{Core/ArchitectureObstruction} & Stable re-export names for fold pattern \\
%
\texttt{Core/SemanticTypeObstruction} & \texttt{SPEC\_020}: semantic typing + three instances + corollaries \\
\texttt{SemanticType/*} & Core bundle; depth / adjudication / RI diagonal; hull tags (seed-local exclusion); pair tags; bipartition bridge; simulation pullback + section (\emph{realization}) \texttt{iff} on reachability; complementary tags; antichain; computability scope note \\
\midrule
\texttt{Reachability/InternalOps} & \texttt{ForwardClosed}, \texttt{PreservedBy}, monotone \texttt{ReflTransGen}, relation extension \\
\texttt{Reachability/ClosureHull} & \texttt{reachableFrom}, hull algebra (union, intersection, idempotence), subrelation monotonicity \\
\texttt{Reachability/Invariants} & Hull--invariant wrappers \\
\texttt{Reachability/ReflectiveSteps} & \texttt{reflectiveSlotStep}, sort preservation \\
\texttt{Reachability/ReflectiveCalculus} & Rich calculus, admissible unions, sub-hull containment \\
\texttt{Reachability/GeneratedCalculi} & \texttt{unionGen}, reflective Bool generators \\
\texttt{Reachability/ArchitectureClasses} & \texttt{Architecture}, \texttt{FoldObstructionBundle} \\
\midrule
\texttt{Obstruction/Fold} & Summit: fold obstruction, hull exclusion, relation-extension persistence \\
\texttt{Obstruction/CanonicalInstances} & Boundary, connectivity, parity, holonomy, sort corollaries \\
\texttt{Obstruction/ReflectiveFold} & Sort mismatch flagship, tower unboundedness \\
\texttt{Obstruction/Architecture\-Universality} & Bundled fold + hull theorems \\
\texttt{Obstruction/OpenCompactWitness} & Discrete witness-step obstruction on $\mathbb{N}$ \\
\midrule
\texttt{Invariants/*} & Sort separation, transport, boundary type, connectivity, orientability-like, homeomorph transport \\
\texttt{Topology/*} & Models, local 1D/2D, holonomy tags, holonomy phase, Hausdorff \\
\midrule
\texttt{Examples/Representational\-Incompleteness} & RI bridge: \texttt{toReflectiveSystem}, sort fold \\
\texttt{Examples/CylinderMobius} & Holonomy parity instance \\
\texttt{Examples/ObserverBridge} & Schematic RFO $\to$ OE route \\
\bottomrule
\end{tabular}
\end{center}
\normalsize

\noindent
All modules compile with zero \texttt{sorry} under
\texttt{lake build ReflectiveFoldObstruction}.

\end{document}
